% 本文件由 LaTeX 排版 Agent 根据有效 translation.json 自动生成。
% author=John Chrysostom; work_id=1918; title=与教宗依诺增爵一世通信集

\chapter{\WorkTitle{与教宗依诺增爵一世通信集}}

% structure: p0001 source=h1 target=omitted reason=page-title-already-rendered-by-parent

% structure: p0002 source=h2 target=section reason=relative-heading-rank; base=h2

\section{\Person{若望}致\Person{依诺增爵}}

% structure: p0003 source=h3 target=subsection reason=relative-heading-rank; base=h2

\subsection{致我主，至为可敬、蒙天主眷爱的主教\Person{依诺增爵}，\Person{若望}在主内致候。}

1. \Emph{我想}，甚至在我们去信之前，尊驾已听说此处所行的不义。因为我们患难之重大，几乎未使世上任何一个角落不知晓此悲惨悲剧：因为报告将所发生之事传至地极，到处引起巨大哀恸与悲叹。但既然我们不应悲哀，而应恢复秩序，并明了藉何种方法才可止住教会这场最惨烈的风暴，我们以为有必要说服我主，即最受尊敬且虔诚的主教\Person{德默特琉}、\Person{盘索斐乌}、\Person{巴普斯}和\Person{欧基纽}，离开他们自己的教会，冒此大海航行之险，离家远行，急速前往仁爱那里，并在清楚告知一切后，尽快采取措施补救这些祸患。我们派了最受尊敬且亲爱的执事\Person{保禄}和\Person{齐利亚古}与他们同去，但我们自己也藉书信之形式，将所发生之事简要告知仁爱。因为\Person{泰奥菲禄}被托付管理亚历山大教会，奉命独自前往君士坦丁堡，有些人向至为虔诚的皇帝控告他，他到达时带着不少埃及主教，仿佛从一开始就表明他是为战争和对抗而来；而且，当他踏进伟大而蒙天主眷爱的君士坦丁堡时，他没有按照自古以来通行的习惯和法律进入教会，他没有与我们交往，也不让我们参与他的谈话、祈祷或共融：但他一下船，便匆匆经过教堂的门厅，离开并在城外某处住宿，尽管我们恳切地邀请他和与他同来的人作客（因为一切都已预备好，住所也已提供，一切合宜之物），但他们和他都不答应。我们见此甚为困惑，无法发现这不义敌意的原因；然而我们尽了自己的本分，做了应作之事，不断恳求他见我们，并说出为何一开始就冒如此大的冲突，使全城陷入混乱。但他不愿说明原因，而那些控告他的人又紧迫，我们至为虔诚的皇帝便召见我们，命令我们到城外泰奥菲禄暂居之处，听取对他的指控。因为他们控告他攻击、杀戮及无数其他罪行；但我们知道教父的法律，尊重并敬重此人，又有他本人的信件，证明诉讼不应超出边界，而各省的事务应在省界内处理，因此我们不愿接受审判官之职，而是极力推辞。但他仿佛要加重先前的侮辱，召见我的总执事，以越权之行为，仿佛教会已孀居无主教，藉此人引诱所有圣职人员归向他；各教会变得无人照顾，因为各处的圣职人员逐渐被撤走，被指示提交反对我们的请愿书，并被训练准备控告。做完这些后，他派人传唤我们受审，尽管他尚未洗清对自己提出的指控，此行为直接违反教规及所有法律。\par

2. 但我们知道自己不是被传唤受审（否则我们会多次出庭），而是被带到敌人和对手面前，这由前后发生的一切清楚证明，于是派遣了几位主教到他那里：\Place{佩西努斯}的\Person{德默特琉}、\Place{阿帕美亚}的\Person{欧里修}、\Place{阿庇亚里亚}的\Person{路比基诺}，以及司铎\Person{日尔曼诺}和\Person{塞维鲁}。他们以适合我们的温和态度回答说，我们不拒绝受审，但拒绝出现在公开的敌人和明显的对手面前。因为，一个人尚未收到任何针对我的起诉状，且从一开始就照上述方式行事，自绝于教会、共融与祈祷，训练控告者，引诱圣职人员，使教会荒凉，他又怎能公正地登上审判席呢？审判席对他来说一点也不合适。因为属于埃及的人，不适宜审判那些在色雷斯的人，更何况此人自己正受指控，又是敌人和对手。然而他毫不羞愧，急欲完成他的计划，尽管我们声明愿意在一百甚至一千位主教面前洗清指控，证明自己是无辜的（事实也如此），他仍不同意：但当我们缺席时，当我们正上诉一个主教会议，要求审判，不畏惧听审我们的案件，而只畏惧公开敌意时，他既接受了我们的控告者，又赦免了被我绝罚的人，并从这些尚未洗清自己所犯罪过的人那里接受对我的控诉，并记录了会议记录——所有这些事都违反法律和教规的秩序。但何必长篇大论？他没有停止策划和行事，直到尽显专横之权与权威，将我们逐出城和教会，那时已是深夜，全体民众紧随我们。我被公控官拖着穿过城中心，被强制拖到海边，被推上船，连夜航行，因为我上诉一个主教会议，以求公正审理我的案件。谁听了这些事不流泪呢？即使他有铁石心肠。\par

但是，如我先前所说，我们不但应为已发生的恶事哀叹，更应加以纠正，所以我恳求你的爱德奋起、怜悯，竭尽全力，使祸害在此处终止。因为即使在我所提及的那些事之后，他仍没有停止他的不法行为，反而企图重新发动先前的攻击。原来，当最虔诚的皇帝赶走了那些无耻冲入教会的人之后，在场的许多主教因见他们的不义而退入自己的教区，如同逃避吞噬万物的烈火一样逃避这些人的侵袭；我们又被邀请回到城里，回到我们曾遭受不公驱逐的教会，由三十多位主教陪同，我们最虔诚的皇帝也为此派遣了一名书记官，而\Person{提阿非罗}立刻逃遁。他出于何种目的、何种原因？我们进城后，恳求最虔诚的皇帝召集一次主教会议，以追究近期事件中的犯法者。因此，提阿非罗对自己所做的事心知肚明，惧怕被定罪，便在皇帝谕旨已发往各地、召集众人从四面八方会集之际，于深夜偷偷跳上一只小船，带着他所有同伙逃走了。\par

3. 但即使在那时，我们仍没有停止，因我们良心无愧，再次向最虔诚的皇帝提出同样的请求；而他也依其虔诚行事，再次派遣使者到埃及，召\Person{提阿非罗}及其同伙前来，就近期事件作出交代，并告知他不要以为他在我们缺席、并违反如此多教会法规的情况下不公正地单方面行径足以成为他的辩护。然而，他并未服从皇命，而是留在家里，以民众叛乱为借口，又推说某些依附于他的人不合时宜的热心（如他所声称的）；然而，在皇帝谕旨到达之前，这些民众早已对他唾骂不已。但我们此刻并不太计较这些事，之所以说这些，只是希望证明他确实被阻止了其恶行。然而，即使在此之后，我们仍未休息，而是迫切要求设立一个审查询问与辩护的法庭；因为我们说，我们准备证明自己清白无辜，而他们则是公然违法。因为当时与他同来的叙利亚人中，有一些留在了此地；我们上前向他们表示愿意为自己辩护，并多次为此事恳求他们，要求将（近期事件的）记录交给我们，或者公布正式的起诉书、指控的性质，或者指出控告者本人；然而，我们未获得其中任何一项，却再次被逐出教会。我该如何讲述随后发生的事呢？它们超越了一切悲剧！什么样的言语能够描述这些事件？什么样的耳朵能不颤栗地聆听？因为当我们坚持这些要求时（如前所述），一大群士兵在伟大的安息日本身，在一天将尽、黄昏临近之时，闯入教堂，暴力驱逐了所有与我们在一起的神职人员，并用武器包围了圣所。从祈祷室出来的妇女，当时刚脱衣准备领受圣洗，因这凶猛的袭击而惊恐地赤身逃跑，不允许穿上合乎妇女的端庄衣饰；事实上，许多人在被驱逐前已受伤，圣洗池中满是鲜血，圣水也被染红。苦难并未在此止息；士兵中，据我们所知，有些尚未领洗，他们进入存放圣器的地方，看见了里面的一切；在如此混乱中，基督至圣的血（正如可能发生的那样）溅在了上述士兵的衣服上；各种暴行如同在野蛮人的围攻中一样发生。平民被赶到旷野，众人都停留在城外，在这盛大的庆节中，教堂变得空荡荡；与我们一起的四十多位主教，无辜地被与民众和神职人员一同驱逐。到处是哭号与哀叹，泪如泉涌，在集市、房屋、荒野之地，城的每一角落都充满了这些灾祸；因这暴行过度，不仅受害者，连那些未遭受任何类似遭遇的人也同情我们，不仅与我们持相同见解的人，连异端者、犹太人和希腊人也是如此；所有地方都陷入骚动、混乱与哀哭，仿佛城已被武力攻占。这些事是违背我们最虔诚的皇帝意愿的，在夜色掩护下发生的；主教们策划了这一切，在许多地方亲自指挥攻击，他们毫不羞愧地让警官代替执事走在他们前面。天亮时，全城的人都迁移到城墙外，在树下和树林中庆祝庆节，如同四散的羊群。\par

4. 此后发生的一切，我让你去想象；因为如前所述，不可能逐一描述每个细节。最糟糕的是，这些巨大而严重的祸患，至今仍未得到遏止，也看不到遏止的希望；相反，祸害每天都在蔓延，我们成了众人的笑柄，或者更确切地说，即使是最无法无天的人也没有人发笑，而是像我所说的，所有人都为这种无法无天的新花样——我们一切苦难的终结——而哀悼。\par

对于其他教会中的混乱，还有什么可说的呢？因为邪恶并未止于此，反而蔓延至东方。正如某种有害体液从头部排出时，身体其余部分也会败坏；同样，这些邪恶如今也起源于这座大城市，仿佛从泉源涌出，混乱已扩散到各处。圣职人员处处起来反抗主教，主教与主教之间、民众与民众之间出现了分裂，将来还会更多。每个地方都在遭受灾祸的剧痛，整个文明世界颠倒了。最可敬、最虔诚的诸位主上，得知这一切后，请展现你们应有的勇气与热忱，以制止这场施加于教会的不法大攻击。因为如果这种风气盛行，并且允许任何人随意进入如此遥远的异教区，驱逐他们想驱逐的人，并按其专断权力为所欲为，那么请放心，一切将毁灭，一种无情的战争将席卷全世界，人人互相攻击，又轮番被攻击。因此，为防止这样的混乱笼罩整个大地，请应允我们的恳求，请以书面声明：这些在我们缺席、仅听一面之词后（尽管我们并未拒绝审判）执行的非法行为是无效的（实际上按其性质就是无效的），而那些犯下如此不义之事的人必须受教会法律的惩罚；至于我们，既未被发现、未被定罪，也未被证明应受处罚，愿我们继续享有你们的通信、你们的爱情以及我们以往所享有的一切。但即使如今，那些犯下如此非法行为的人愿意披露指控的依据（他们据此不义地驱逐了我们），而既未提供备忘录，也未提供正式起诉书，控告者亦未出庭；然而，若能成立一个公正的法庭，我们愿意接受审判，进行辩护，并证明自己清白，正如我们确实是清白的。因为他们所做的一切都超出了任何秩序、任何教会法律和教规的界限。我为何提教会教规？即使在外教的法庭上，也从未犯下如此胆大妄为的行为，甚至蛮族的法庭也没有。无论是\Place{西徐亚人}还是\Place{萨尔马提亚人}，都绝不会以这种方式审理案件：仅听一面之词，在被告缺席的情况下作出判决，而被告只是拒绝敌意，并非拒绝案件审理，并愿意召集任何数量的法官，声称自己无辜，并能在世人面前洗清指控，证明自己在各方面都清白无过。\par

因此，思考这一切，并由我主最虔诚的弟兄主教们清楚得知所有详情后，愿您被促使为我们而发挥您的热忱。因为如此行，您不仅恩赐于我们，也恩赐于整个教会，并必从为使教会和平而行事的天主那里获得赏报。祝您永远安好，最可敬、最神圣的主上，请为我祈祷。\par

% structure: p0011 source=h3 target=subsection reason=relative-heading-rank; base=h2

\subsection{致罗马主教\Person{依诺增爵}，在主内问安}

我们的身体固然停留在一处，但爱情的翅膀却飞遍世界每一个角落。同样，我们虽在如此漫长的旅途上分离，但仍亲近你的陛下，与你日日共融，以爱情的目光注视你灵魂的勇气、你品格的纯正、你的坚定与不屈，以及你不断赐予我们的恒常不变的极大安慰。因为浪涛越高，暗礁越多，暴风越众，你的警醒也越强：我们之间的漫长旅途、大量时间的消耗、处理事件的困难，都未能使你变得懈怠；你继续效法最优秀的舵手，他们在看见浪花翻腾、海面汹涌、水势猛烈、白昼如黑夜时，最是警觉。因此我们对你满怀感激，渴望向你寄送如雨般的书信，以此给予我们最大的满足。但由于此地荒凉，我们无法做到这一点；（因为不仅从你们地区来的人，就连住在我们这世界部分的人，也无人能轻易与我们交往，既因为距离遥远——我们被囚之地位于国家最偏僻的角落——又因盗匪的恐怖阻挡了整个旅程：）我们恳求你因我们长久的沉默而怜悯我们，而不是因此谴责我们懒惰。为证明我们的沉默并非出于疏忽，我们在长时间后终于找到了我们最尊敬和亲爱的\Person{若望}长老和\Person{保禄}执事，并通过他们致信，继续向你表达感激，因为你在对我们的善意和热忱上甚至超越了慈爱的父母。确实，就你的陛下而言，一切本应得到适当修正，邪恶和过犯的积累本应被清除，教会本应享有和平与镜面般的平静，万事本应顺流而行，被藐视的法律和被违反的教父法令本应得到辩护。但既然实际上这些事一件也没有发生，那些先前作恶的人努力加重从前的不义，我略过对他们后续行为的详细叙述：因为叙述不仅会超出信函的限度，甚至超出历史的限度；只有这一点我恳求你警醒的灵魂：即使那些使一切混乱的人顽固不化、不可救药地败坏，也不要让那些承担医治他们责任的人变得灰心气馁，当他们想到要完成之事的重大时。因为你眼前的这场争战几乎要为全世界而战，为被夷为平地的教会、四散的人民、受攻击的圣职人员、被流放的主教、被违反的祖传法律而战。因此我们恳求你的勤勉，一次、两次，甚至多次，随着风暴加剧，要表现出更大的热忱。因为我们期待为修正这些错误会做更多的事。但即使这没有发生，你至少已拥有仁慈的天主为你预备的冠冕，你的爱情所给予的抵抗对受冤屈者将是莫大的安慰：因为现在我们已在流放中度过第三年，遭受饥饿、瘟疫、战争、持续围攻、难以言表的孤独、每日的死亡和\Place{伊索里亚}人的刀剑，我们因你性情和信赖的恒常不变，以及沉浸在你丰沛而真诚的爱情中，得到不少鼓励和安慰。这就是我们的防护墙，我们的保障，我们平静的港湾，我们无尽福佑的宝藏，我们的喜乐和极大欢乐的缘由。即使我们再次被带到比这里更荒凉的地方，我们也会将这份爱情随身带去，作为我们受苦的不小安慰。\par

% structure: p0013 source=h2 target=section reason=relative-heading-rank; base=h2

\section{从\Person{依诺森}致\Person{若望}}

% structure: p0014 source=h3 target=subsection reason=relative-heading-rank; base=h2

\subsection{致爱弟\Person{若望}，\Person{依诺森}}

虽然无辜的人理当期待一切美好事物，并切望天主的仁慈，但我们也劝勉顺从，由执事\Person{齐里亚古}之手送出一封合宜的信函；好使傲慢在压迫上的力量，不大于善良良心在持守希望的力量。因为你是如此众多人民的教师和牧者，无需他人教导：最优秀的人常常受到考验，看他们是否能在完美的忍耐中坚持，不被任何痛苦劳碌所征服；确实，良心是对付一切不义临到我们的坚强防御；除非一个人以忍耐克服这些，他便为恶意的猜疑提供了论据。因为那首先信赖天主、其次信赖自己良心的人，应当忍受一切；既然高尚良善的人能够特别被训练到忍耐，因为圣经护卫他的心智；我们向人民传授的神圣课程充满例证，证实几乎所有圣人都在不同方式下不断受压迫，并好似经过某种检验而受考验，如此达到忍耐的冠冕。让良心本身安慰你的爱情，最可敬的兄弟，因为良心在苦难中提供德行的安慰。因为在主基督的眼下，良心被净化后，将在平安的港湾中找到安息。\par

% structure: p0016 source=h3 target=subsection reason=relative-heading-rank; base=h2

\subsection{\Person{依诺森}，主教，致长老与执事，并致\Place{君士坦丁堡}教会全体圣职人员和人民，即服从主教\Person{若望}的可敬弟兄们，请安。}

您由\Person{哲曼努斯}司铎和\Person{卡西亚努斯}执事之手所寄送的信函，使我以焦虑关切的心情研读了您展现在我眼前的灾难景象；通过反复阅读您的描述，我彻底领悟到您的信德正如何处于巨大的痛苦与劳苦之中：此事唯有以忍耐的安慰才能治愈；因为我们的天主将迅速终结如此大的苦难，并助您承受这一切。此外，在赞扬您的案情陈述中包含许多鼓励忍耐的证言时，我注意到在您仁爱的书信开头便放置了这一必要的安慰：因为我们本应写信安慰您，您却以您的来信抢先了。因为我们的主惯于将这种忍耐赐予受苦之人，使基督的仆役在苦难中能藉着反思自己所受的苦也发生于昔日的圣人们身上而安慰自己。我们也必能从您的书信中获得安慰：因为我们并非与您隔绝而不表同情，因为我们也在您们身上受惩罚。因为谁能忍受那些本应特别热心促进教会安宁及自身和谐之人所犯的罪行呢？如今，由于习俗颠倒，无辜的司铎从他们自己教会的领导位置上被驱逐。这正是您们的首席弟兄、同工、您们的主教\Person{若望}所不公正地遭受的，他没有得到任何听审：没有罪行指控他，也没有听到任何指控。这不义计谋的目的是什么？为了避免可能或寻求审判的借口，竟将他人引入活着的司铎的位置，仿佛那些始于此类罪行的人能被人认为拥有任何善或做出任何正确的事。因为我们明白，此类行为从未由我们的先辈实施过；更确切地说，它们之所以被阻止，是因为没有人被授权祝圣另一人去取代仍然活着的人。因为一个虚假的祝圣不能剥夺司铎的品秩：因为被错误代替他人的人也不能成为主教。关于教规的遵守，我们声明应当遵循那些在\Place{尼西亚}所制定的教规，唯独这些教规是\OriginalTerm{天主教会}{Catholica Ecclesia}必须服从和承认的。如果有人提出其他与尼西亚所制定的教规相悖并且被证明由异端者制定的法规，天主教会的主教们应当拒绝它们。因为异端者的发明不应附加在天主教教规上；因为他们总是以敌对和非法的法令意图削弱尼西亚教规的设计。因此我们不仅说这些不应被遵循，而且应当像从前在\Place{萨迪卡}会议中由我们之前的主教们所做的那样，将它们归入异端和分裂的法令中予以谴责。最受尊敬的弟兄们，与其让直接违背教规的行为具有任何效力，不如谴责善行更为合适。但如今面对这些事我们该怎么办？需要召开一次主教会议来裁决它们，我们早已宣布应当召开会议，因为这是平息此类风暴骚动的唯一方法：如果我们得到这个会议，则将这些邪恶的治愈托付于伟大的天主及祂的基督、我们的主之旨意是合适的。因魔鬼嫉妒为使信徒受考验而引起的所有骚动都将平息；通过我们信德的坚定，我们不应在任何事上对主失望。因为我们也正在积极思考如何召开普世主教会议，以便藉天主的旨意终结这些扰乱运动。因此让我们忍耐一时，以忍耐的壁垒巩固自己，期望一切能藉我们天主的助佑恢复给我们。此外，您所说所经历的一切，我们从已逃往\Place{罗马}的同道主教们那里通过确切询问得知，虽然大部分时间不同，即\Person{德默特琉}、\Person{齐里亚库斯}、\Person{尤利西乌斯}和\Person{帕拉迪乌斯}，他们现与我们同在。\par

\SourceRecord{Translated by W.R.W. Stephens.；Nicene and Post-Nicene Fathers, First Series；Vol. 9.；Edited by Philip Schaff.；Buffalo, NY: Christian Literature Publishing Co.,；1889.；<http://www.newadvent.org/fathers/1918.htm>.}
